% Zone „Das kennst du schon“ – aus bank/rationale-zahlen/zone.jsonl (zusammenbau v0.1)
\einheitenkopf[zone][Kennst du schon]{Das kennst du schon}
\setcounter{aufgabe}{0}

%% TODO Titel: Ich-kann-Satz fehlt in der Bank (Platzhalter: Kettenname) – Nr. 1
%% TODO Anweisung: ein Satz über der ersten Teilaufgabe fehlt in der Bank – Nr. 1
\begin{aufgabe}{Natürliche Zahlen addieren, subtrahieren, multiplizieren, dividieren im Kopf (Einmaleins)}
\begin{teile}
\teil Berechne: $7 \cdot 8$ \leerfeld
\teil Berechne: $120 : 8$ \leerfeld
\end{teile}
\end{aufgabe}

%% TODO Titel: Ich-kann-Satz fehlt in der Bank (Platzhalter: Kettenname) – Nr. 2
%% TODO Anweisung: ein Satz über der ersten Teilaufgabe fehlt in der Bank – Nr. 2
\begin{aufgabe}{Zahlenstrahl}
\begin{teile}
\teil Welche Zahl gehört zu A? Lies ab.

\zahlenstrahl[xmin=0,xmax=10,xstep=1]{7/A}
A = \leerfeld
\teil Welche Zahl gehört zu C? Lies ab.

\zahlenstrahl[xmin=0,xmax=2,xstep=0.1,karo=0.5]{1.3/C}
C = \leerfeld
\end{teile}
\end{aufgabe}

%% TODO Titel: Ich-kann-Satz fehlt in der Bank (Platzhalter: Kettenname) – Nr. 3
%% TODO Anweisung: ein Satz über der ersten Teilaufgabe fehlt in der Bank – Nr. 3
\begin{aufgabe}{Dezimalzahlen und Brüche addieren, subtrahieren, multiplizieren (Kommazahlen addieren, Bruch mal ganze Zahl)}
\begin{teile}
\teil Berechne: $2{,}5 + 1{,}3$ \leerfeld
\teil Berechne: $4{,}6 - 1{,}85$ \leerfeld
\end{teile}
\end{aufgabe}

%% TODO Titel: Ich-kann-Satz fehlt in der Bank (Platzhalter: Kettenname) – Nr. 4
%% TODO Anweisung: ein Satz über der ersten Teilaufgabe fehlt in der Bank – Nr. 4
\begin{aufgabe}{Punkt vor Strich und Klammern mit natürlichen Zahlen}
\begin{teile}
\teil Berechne: $3 + 4 \cdot 5$ \leerfeld
\teil Berechne: $0{,}5 + 4 \cdot 2{,}5$ \leerfeld
\end{teile}
\end{aufgabe}

%% TODO Titel: Ich-kann-Satz fehlt in der Bank (Platzhalter: Kettenname) – Nr. 5
%% TODO Anweisung: ein Satz über der ersten Teilaufgabe fehlt in der Bank – Nr. 5
\begin{aufgabe}{Wert eines Terms mit Platzhalter berechnen (Vorzahl mal x plus Zahl, für einen gegebenen Einsetzwert)}
\begin{teile}
\teil $3 \cdot x + 2$ für $x = 4$ – welcher Wert? \leerfeld
\teil $2 \cdot x + 1$ für $x = 2{,}5$ – welcher Wert? \leerfeld
\end{teile}
\end{aufgabe}

%% TODO Titel: Ich-kann-Satz fehlt in der Bank (Platzhalter: Kettenname) – Nr. 6
%% TODO Anweisung: ein Satz über der ersten Teilaufgabe fehlt in der Bank – Nr. 6
\begin{aufgabe}{Punkt vor Strich und Klammern mit natürlichen Zahlen – Fehler finden}
\begin{teile}
\teil Mia rechnet: $27 - 9 : 3 = 18 : 3 = 6$. Finde den Fehler.
\end{teile}
\end{aufgabe}

%% TODO Titel: Ich-kann-Satz fehlt in der Bank (Platzhalter: Kettenname) – Nr. 7
%% TODO Anweisung: ein Satz über der ersten Teilaufgabe fehlt in der Bank – Nr. 7
\begin{aufgabe}{Punkt vor Strich und Klammern mit natürlichen Zahlen – selbst rechnen}
\begin{teile}
\teil Berechne: $40 - 18 : 6$ \leerfeld
\end{teile}
\end{aufgabe}
